\section{Modellierungssprachen}
In der Informatik nutzen wir Modellierungssprachen um die Anforderungen an ein Softwaresystem, als auch dessen innere Abläufe und Strukturen in einer höheren Abstraktionsebene festzulegen. Modellierungssprachen finden Anwendung in der
\begin{itemize}
	\item \textbf{Entwurfsphase}, wenn die Anforderungen ermittelt und festgelegt werden
	\item \textbf{Implementierungsphase}, wenn die Softwareentwickler/innen die Software in einer höheren Programmiersprache umsetzen
	\item \textbf{Dokumentation} für Kunden oder Kollegen (weil sie ja eh schon angefertigt wurden, kann man diese hier prima weiter verwenden)
\end{itemize}
% ------------------------------------------------------------------------------------------
\subsection{Unified Modeling Language}
Die \ac{UML} ist eine grafische Modellierungssprache zur Spezifikation, Konstruktion, Dokumentation und Visualisierung von Software-Teilen und anderen Systemen. Sie ist derzeit in der Major-Version 2 in der ISO 19505 spezifiziert. In Abb. \ref{fig:uml_diagrammhierarchie} ist eine Übersicht über die gängigsten UML-Diagramme dargestellt. UML gliedert sich im Wesentlichen in 
\begin{itemize}
	\item \textit{Strukturdiagramme}, um die Struktur (z. B. Klassen, Komponenten) einer Software zu beschreiben
	\item \textit{Verhaltensdiagramme}, um das Verhalten einer Software (z. B. Aktivitäten oder Zustände) zu beschreiben
	\item \textit{Interaktionsdiagramme}, um die Interaktion zwischen Software-Komponenten, aber auch zwischen Nutzern zu modellieren
\end{itemize}

\begin{figure}[H]
	\centering
	\includegraphics[width=1\textwidth]{\imgpath/uml_diagrammhierarchie.png}
	\caption{UML-Diagrammhierarchie (Wikipedia)}
	\label{fig:uml_diagrammhierarchie}
\end{figure}

% ------------------------------------------------------------------------------------------

\subsection{Draw.io einrichten}

\begin{tipbox}{Software}
	\textbf{Bevor du beginnst:} Nutze zur Modellierung die Online-Software \textbf{draw.io}. Blende zuvor die UML 2.5 Symbole ein und stelle sicher, dass du im Diagramm die \textbf{richtigen Symbole} verwendest!	
\end{tipbox}

\begin{importantbox}{Symbole}
	Wir sind Profis. Die Nutzung der richtigen Symbole gehört genauso zu unserer Professionalität wie die fachliche/inhaltliche Richtigkeit unserer Arbeit.	
\end{importantbox}

\begin{figure}[H]
	\centering
	\includegraphics[width=1\textwidth]{\imgpath/drawio_shapes.jpg}
	\caption{Draw.io UML 2.5 Symbole aktivieren}
	\label{fig:drawio_shapes}
\end{figure}
